%==============================================================================
%  TAD Publications / TRB 论文格式  —— LaTeX 模板
%  本模板复刻自 Word 样例 (TAD_Publications_AMSamplePaper.docx)
%  编译方式：pdflatex paper.tex   (运行两次以更新页码/交叉引用)
%
%  你只需要在下面带 ▶ 标记的地方替换成自己的内容即可。
%==============================================================================
\documentclass[12pt,letterpaper]{article}

%--- 字体：Times New Roman 等价字体（正文与公式都用 Times 风格）-----------------
\usepackage{mathptmx}
\usepackage[T1]{fontenc}
\usepackage[utf8]{inputenc}

%--- 页面：US Letter，四周 1 英寸页边距，页眉/页脚 0.5 英寸 ----------------------
\usepackage[letterpaper,margin=1in,headheight=15pt,headsep=0.25in,footskip=0.3in]{geometry}

%--- 页眉（斜体 running head）+ 页脚（居中页码）---------------------------------
\usepackage{fancyhdr}
\pagestyle{fancy}
\fancyhf{}
\fancyhead[L]{\itshape Author, Author, and Author}   % ▶ 改成你的 running head
\fancyfoot[C]{\thepage}
\renewcommand{\headrulewidth}{0pt}
\renewcommand{\footrulewidth}{0pt}

%--- 正文：左对齐(右侧自然参差)、单倍行距、首行缩进 0.5 英寸、段间不额外加空 --------
\usepackage{ragged2e}
\setlength{\RaggedRightParindent}{0.5in}
\RaggedRight
\setlength{\parindent}{0.5in}
\setlength{\parskip}{0pt}
% 如果你更喜欢两端对齐(齐边)的“正式出版”观感，把上面的 \RaggedRight 注释掉即可。

%--- 标题层级样式 --------------------------------------------------------------
\usepackage{titlesec}
\setcounter{secnumdepth}{0}                                   % 标题不编号
% 一级标题：全部大写 + 加粗，顶左对齐
\titleformat{\section}{\normalfont\bfseries}{}{0pt}{\MakeUppercase}
\titlespacing*{\section}{0pt}{12pt}{4pt}
% 二级标题：加粗
\titleformat{\subsection}{\normalfont\bfseries}{}{0pt}{}
\titlespacing*{\subsection}{0pt}{12pt}{4pt}
% 三级标题：斜体
\titleformat{\subsubsection}{\normalfont\itshape}{}{0pt}{}
\titlespacing*{\subsubsection}{0pt}{12pt}{4pt}

%--- 图表标题：标签与正文一起加粗，空格分隔，左对齐 ------------------------------
\usepackage{graphicx}
\usepackage{booktabs}
\usepackage{array}
\usepackage{caption}
\captionsetup{labelsep=space,font=bf,justification=raggedright,singlelinecheck=false}
\captionsetup[table]{name=TABLE}     % 表标题前缀：TABLE 1 …（在表上方）
\captionsetup[figure]{name=Figure}   % 图标题前缀：Figure 1 …（在图下方）

%--- 超链接 --------------------------------------------------------------------
\usepackage[hidelinks]{hyperref}

%--- 自定义命令 ----------------------------------------------------------------
% 结构化摘要的加粗引导词
\newcommand{\abshead}[1]{\textbf{#1}\ }
% 参考文献：悬挂缩进（首行顶格，回行缩进 0.5 英寸）
\newcommand{\refitem}[1]{\par\noindent\hangindent=0.5in\hangafter=1 #1\par\vspace{6pt}}

%==============================================================================
\begin{document}

%------------------------------------------------------------------------------
%  题名与作者信息（全部左对齐，单倍行距）
%------------------------------------------------------------------------------
\begin{flushleft}
{\bfseries Title of TRB Paper Format Example}\par   % ▶ 论文标题
\vspace{12pt}

% ▶ 作者块 1：学术界作者
{\bfseries Academic Author Name\textsuperscript{*}}\par
Department of XXX\par
Institution Name, City, State or Country, and Postcode\par
ORCID: https://XXX\par
Email: author@university.edu\par
\vspace{12pt}

% ▶ 作者块 2：公共部门作者
{\bfseries Public Sector Author Name}\par
Position\par
XXStateXX Department of Transportation\par
Department (if applicable)\par
City, State or Country, Postcode\par
Email: abc@dot.gov\par
\vspace{12pt}

% ▶ 作者块 3：业界从业者
{\bfseries Private Practitioner Author Name}\par
Position\par
Company\par
City, State or Country, Postcode\par
Email: enj@abc.com\par
\vspace{12pt}

Total Number of Pages: XX\par                       % ▶ 总页数
\vspace{12pt}

\textit{Submitted [Submission Date]}\par             % ▶ 提交日期
\textsuperscript{*}Corresponding Author\par
\end{flushleft}

\vspace{6pt}

%------------------------------------------------------------------------------
%  摘要（结构化摘要，两端对齐）
%------------------------------------------------------------------------------
\section{Abstract}
\noindent\textit{(Must be formatted as a structured abstract per the instructions
for authors. Must be less than 300 words. The content below is for illustrative
purposes only.)}

\justifying  % 摘要内部两端对齐

\abshead{Objectives:}Fare capping, a policy in which a transit agency caps the
maximum amount a rider pays over a given period, has emerged as a recent innovation
in public transit fare policy. However, its use across U.S.\ transit agencies and
its rider benefits are not yet well documented. The objective of this research is to
compile and compare fare capping policies across the U.S.\ and to explore the
potential rider benefits.

\abshead{Methods:}A multiple case study method was used to examine fare capping
policies at the 101 largest transit agencies in the U.S. Agency fare policies were
identified through publicly available documents on the internet. For each agency,
fare policy documents were reviewed to identify whether daily, weekly, and monthly
caps were offered, and the number of one-way adult fares required to reach each cap
were calculated.

\abshead{Findings:}As of September 2021, fare capping was adopted by at least 21 of
the 101 largest U.S.\ transit agencies. Of those agencies, twenty used daily caps,
four used weekly caps, and fourteen used monthly caps. The results also revealed that
most daily fare caps may not offer much, if any, discount for a typical commuter
assumed to make only two transit trips per day; however, daily fare capping likely
provided important discounts to transit-dependent riders, tourists, and other
frequent riders. Some innovative fare capping policies like ``nested'' fare capping,
which refers to fare caps within fare caps, were also identified.

\abshead{Novelty:}This paper provides one of the first systematic syntheses of fare
capping policies across major U.S.\ transit agencies. It also introduces a
comparative framework based on cap structure, trip thresholds, and rider discount
levels.

\abshead{Practical Applications:}The findings can help transit agencies design fare
capping policies that better match rider travel patterns and policy goals.

\RaggedRight  % 摘要结束，正文恢复左对齐

%------------------------------------------------------------------------------
%  正文
%------------------------------------------------------------------------------
\section{Introduction}
Sed efficitur magna nisl, id aliquet leo molestie et. Quisque vestibulum imperdiet
neque, id cursus felis rutrum et. Duis porttitor nisl arcu, sit amet varius justo
convallis vitae. Maecenas pretium ullamcorper massa sed volutpat (Smith 2020). Duis
libero augue, porttitor sed magna quis, ultrices molestie magna. Nam in orci in dolor
accumsan rhoncus ut et leo. Phasellus rutrum enim ut purus lacinia, eget faucibus
odio fermentum. Vivamus at dolor ante. Suspendisse porttitor diam in diam
condimentum, et commodo diam semper. Pellentesque habitant morbi tristique senectus
et netus et malesuada fames ac turpis egestas.

Morbi dictum placerat erat ac tempor. Pellentesque elementum accumsan quam. Integer
ultricies feugiat placerat. Morbi quis eros et augue iaculis auctor. Aliquam rhoncus
lacus ut tortor egestas, in luctus nisl semper. Quisque ultrices sapien turpis, at
aliquam est sodales sit amet. Nulla congue risus in ligula aliquam cursus. Sed vel
augue consequat orci mollis aliquet nec non arcu. Ut porttitor urna urna. Maecenas
tempor, augue quis gravida suscipit, tortor nulla euismod turpis, at sagittis leo
nisi quis justo. Ut ut elementum velit. Ut nec feugiat urna.

\section{Methods}
Mauris scelerisque maximus viverra. Proin sagittis faucibus viverra. Ut et tellus
eros. Mauris sollicitudin porttitor enim vel congue. Quisque dictum ornare posuere.
Fusce malesuada rhoncus nisl sit amet rutrum. Vivamus posuere elit dui, eu
scelerisque urna interdum sed. Suspendisse potenti. Nulla vel massa non diam
dignissim tristique. Aenean non tempus orci. Aliquam lobortis posuere nisi a
imperdiet.

Lorem ipsum dolor sit amet, consectetur adipiscing elit. Morbi dictum placerat erat
ac tempor. Pellentesque elementum accumsan quam. Integer ultricies feugiat placerat.
Aliquam rhoncus lacus ut tortor egestas, in luctus nisl semper. Quisque ultrices
sapien turpis, at aliquam est sodales sit amet. Maecenas tempor, augue quis gravida
suscipit, tortor nulla euismod turpis, at sagittis leo nisi quis justo.

Duis ut turpis maximus, semper sapien in, elementum metus. Proin sit amet sem sit
amet tortor blandit dapibus. Nam fringilla nunc quis nisl faucibus laoreet. Nam eu
elit laoreet, iaculis sem quis, rhoncus libero (Smith and Jones 2020). Aliquam
maximus nisi risus, a ultricies nunc auctor sit amet. Donec id dui in mi malesuada
volutpat. Donec vel mi id mi congue convallis eget et nibh.

\subsection{Level 2 Header}
Quisque ultrices sapien turpis, at aliquam est sodales sit amet. Nulla congue risus in
ligula aliquam cursus. Sed vel augue consequat orci mollis aliquet nec non arcu. Sed a
neque sodales, rhoncus massa vel, imperdiet augue. Ut porttitor urna urna. Maecenas
tempor, augue quis gravida suscipit, tortor nulla euismod turpis, at sagittis leo nisi
quis justo. Ut ut elementum velit. Ut nec feugiat urna.

\subsubsection{Level 3 Header}
Fusce malesuada rhoncus nisl sit amet rutrum. Vivamus posuere elit dui, eu scelerisque
urna interdum sed (Smith et al.\ 2020). Etiam quis libero id lacus bibendum interdum.
Morbi et pretium lorem (Smith 2021, 25--27). Quisque mollis, enim vitae dapibus
rhoncus, ex ipsum euismod metus, ultrices pharetra est augue et augue. Pellentesque
commodo maximus nulla, in faucibus nisl vestibulum sed (Boback 2020, 45--78). Duis ut
turpis maximus, semper sapien in, elementum metus (\textbf{Figure 1}), amet sem sit
amet tortor blandit.

% ▶ 插图：把下面的占位框换成 \includegraphics[width=...]{你的图片文件}
\begin{figure}[h]
  \centering
  \fbox{\parbox[c][1.6in][c]{0.7\linewidth}{\centering\itshape [Figure placeholder —
  replace with \textbackslash includegraphics]}}
  \caption{Caption for figure}
  \label{fig:1}
\end{figure}

Maecenas pretium ullamcorper massa sed volutpat. Duis libero augue, porttitor sed
magna quis, ultrices molestie magna. Nam in orci in dolor accumsan rhoncus ut et leo.
Phasellus rutrum enim ut purus lacinia, eget faucibus odio fermentum. Vivamus at dolor
ante. Suspendisse porttitor diam in diam condimentum, et commodo diam semper.

\section{Results}
Nulla eu urna et erat bibendum venenatis sit amet sed sem. Nam quis quam tempor,
placerat turpis a, imperdiet ante. Integer eleifend aliquam magna quis viverra. Nulla
consequat sem vitae mi lacinia maximus. Fusce id ullamcorper nibh. Mauris sagittis,
nunc et laoreet condimentum, lectus leo bibendum nisi, eu posuere nulla nisl non
libero. Vestibulum ante ipsum primis in faucibus orci luctus et ultrices posuere
cubilia Curae.

Proin ut lacus in diam aliquet ultrices a eu neque. Fusce nec maximus elit, tempus
ultrices sem. Cras vitae mi commodo, placerat justo a, gravida est. Sed gravida justo
velit, et pharetra sem placerat in. Vivamus laoreet interdum magna eget luctus. Morbi
volutpat pellentesque libero.

Sed eu lacus sit amet risus sodales pretium a ac ex. Donec gravida auctor efficitur.
Curabitur maximus pellentesque mollis. Nam eu ante sit amet ex feugiat dictum eget eu
nisl. Maecenas hendrerit eget sem ut blandit. Nam sit amet dolor volutpat, sollicitudin
mi iaculis, bibendum felis (\textbf{Table 1}).

% ▶ 表格：表标题在表上方
\begin{table}[h]
  \caption{Measurement Conversion}
  \label{tab:1}
  \centering
  \begin{tabular}{lll}
    \toprule
    \textbf{When You Know} & \textbf{Multiply by} & \textbf{To Find} \\
    \midrule
    \textbf{Length} & & \\
    inches (in.)            & 25.4  & millimeters (mm) \\
    feet (ft)               & 0.305 & meters (m) \\
    yards (yd)              & 0.914 & meters (m) \\
    miles (mi)              & 1.61  & kilometers (km) \\
    \addlinespace
    \textbf{Area} & & \\
    square inches (in.\textsuperscript{2}) & 645.1 & millimeters squared (mm\textsuperscript{2}) \\
    square feet (ft\textsuperscript{2})    & 0.093 & meters squared (m\textsuperscript{2}) \\
    square yards (yd\textsuperscript{2})   & 0.836 & meters squared (m\textsuperscript{2}) \\
    acres                                  & 0.405 & hectares (ha) \\
    square miles (mi\textsuperscript{2})   & 2.59  & kilometers squared (km\textsuperscript{2}) \\
    \addlinespace
    \textbf{Volume} & & \\
    fluid ounces (fl oz)              & 29.57 & milliliters (mL) \\
    gallons (gal)                     & 3.785 & liters (L) \\
    cubic feet (ft\textsuperscript{3})  & 0.028 & meters cubed (m\textsuperscript{3}) \\
    cubic yards (yd\textsuperscript{3}) & 0.765 & meters cubed (m\textsuperscript{3}) \\
    \bottomrule
  \end{tabular}
\end{table}

\section{Discussion}
Nunc eros felis, dapibus eu nisl convallis, pellentesque facilisis sapien. Duis vel
magna ac ligula rutrum volutpat. Donec sed malesuada ipsum. Morbi sed metus sem. Etiam
bibendum posuere elementum. In a tempor mauris. Mauris congue in diam ac malesuada.
Praesent sed dignissim erat.

Nunc porttitor magna sed ultrices condimentum. Suspendisse in nisl condimentum, posuere
metus eu, ullamcorper mi. Fusce rutrum ipsum at felis pharetra porttitor. Pellentesque
nec massa nulla \textbf{Equation 1}:

% ▶ 公式：把占位内容替换成你的公式
\begin{equation}
  y = \beta_0 + \beta_1 x_1 + \beta_2 x_2 + \varepsilon
\end{equation}

Sed eu lacus sit amet risus sodales pretium a ac ex. Donec gravida auctor efficitur.
Curabitur maximus pellentesque mollis. Maecenas hendrerit eget sem ut blandit. Nam sit
amet dolor volutpat, sollicitudin mi iaculis, bibendum felis \textbf{Equation 2}:

\begin{equation}
  \hat{p} = \frac{1}{1 + e^{-(\beta_0 + \sum_i \beta_i x_i)}}
\end{equation}

Donec libero velit, bibendum sit amet dignissim in, tempus vel lorem. Donec consectetur
dui lorem, gravida eleifend metus blandit ut. Fusce metus lorem, volutpat sagittis orci
id, lobortis semper lectus. Nunc eu erat porttitor, molestie elit nec, pharetra enim.

\section{Conclusions}
Nulla eu urna et erat bibendum venenatis sit amet sed sem. Nam quis quam tempor,
placerat turpis a, imperdiet ante. Integer eleifend aliquam magna quis viverra. Nulla
consequat sem vitae mi lacinia maximus. Fusce id ullamcorper nibh. Proin id elementum
metus. Cras dolor libero, dignissim in magna vitae, tincidunt porta tortor.

%------------------------------------------------------------------------------
%  文末声明部分
%------------------------------------------------------------------------------
\section{Acknowledgments}
The author must be transparent in their use of AI or LLMs in the Submission Details tab
in Editorial Manager and they must describe which model was used and for what purpose in
their paper's ``Acknowledgements'' section.

\section{Author Contributions}
The authors confirm contribution to the paper as follows: study conception and design:
X.\ Author, Y.\ Author; data collection: Y.\ Author; analysis and interpretation of
results: X.\ Author, Y.\ Author, Z.\ Author; draft manuscript preparation: Y.\ Author,
Z.\ Author. All authors reviewed the results and approved the final version of the
manuscript.

\section{Declaration of Conflicting Interests}
The authors declared no potential conflicts of interest with respect to the research,
authorship, and/or publication of this article.

\section{Funding}
The authors disclosed no financial support for the research, authorship, and/or
publication of this article.

%------------------------------------------------------------------------------
%  参考文献（悬挂缩进；此处手动排版，Chicago 作者-年份样式）
%  若文献多，建议改用 BibTeX/biblatex 管理。
%------------------------------------------------------------------------------
\section{References}
\refitem{Smith, John. 2020. ``Evaluating Pavement Performance Under Variable Climate
Conditions.'' \textit{Journal of Transportation Engineering} 146: 123--135.
https://doi.org/10.0000/xxxxx}

\refitem{Smith, John, and David Jones. 2020. ``Assessing Traffic Flow Improvements Using
Adaptive Signal Control Systems.'' Paper presented at the 99th Annual Meeting of the
Transportation Research Board, Washington, DC.}

\refitem{Smith, John, David Jones, and Thomas Haas. 2020. \textit{Evaluation of Bridge
Deck Deterioration and Maintenance Strategies}. Publication FHWA-RD-20-001. Washington,
DC: Federal Highway Administration, U.S.\ Department of Transportation.}

\refitem{Smith, John. 2021. ``Analysis of Traffic Congestion Patterns in Urban
Corridors.'' \textit{Journal of Transportation Engineering} 147: 25--27.
https://doi.org/10.0000/xxxxx}

\refitem{Boback, Luke. 2020. \textit{Fundamentals of Transportation Systems Analysis}.
Springer. Chapter 3, 45--78.}

\end{document}
